\documentclass{article}
\usepackage{/globals/de}
\begin{document}
Computer können in unterschiedliche Klassen aufgeteilt werden. Im Sinne dieses Kurses wird wie folgt sortiert.
\begin{description}
 \item[Mainframes] Mainframes sind besonders große Computer, mit einer hohen Rechen- und Speicherkapazität. Diese werden in Rechenzentren genutzt.
 \item[Workstations] Workstations sind personal Computer wie man sie im alltag benutzt. Sie haben genug rechenleistung um auch selbst die meisten alltäglichen Anwendungen zu laufen
 \item[Thin-Clients] Thin-Clients, oder auch Netzwerkcomputer (NC) genannt sind eher billigere und leistungsschwächere Computer die ihre Ein- und Ausgaben über ein Netzwerk von an/von einem Server streamen.
 \item[Mobile Endgeräte] Mobile Endgeräte beinhalten alle Rechner die einfach transportiert werden können; Handys, Tables, Laptops, \etc.
 \item[Embedded Systems] Embedded Systems sind kleine, spezialisierte Rechner, die in andere Systeme eingebettet sind, um spezifische Aufgaben zu erledigen. Sie werden oftmals selbst nicht als Rechner wahrgenommen.
\end{description}
\end{document}